% !TeX root = ../../Book5.tex
% 纲要标题：### 18.3 Verma 模
% 纲要位置：工作记录/计划纲要.md，第 4661 行。
% 纲要细目（保留原定写作范围）：
%
% * 诱导构造
% * 唯一简单商
% * 有限维性的支配整性条件
% * 比较 A2 的 PBW 权空间与自然简单商，解释所有真子模共同缺失最高权线的作用
%

\section{Verma 模}
\label{b5:sec:1803}

指定最高权以后，可以只加入最高权向量必须满足的关系，暂不要求表示有限维。由一个同权最高权向量生成的表示都是所得模的商。这样便可先构造最高权模及其简单商，再研究哪些附加关系使商有限维。

\subsection{诱导构造}
\label{b5:subsec:1803:01}

\begin{definition}\label{b5:def:verma-module}
	设 \(\mathfrak g\) 为有限维复半单 Lie 代数，\(\mathfrak h\) 为其 Cartan 子代数，\(\Phi^+\) 为相应根系的一组正根。令 \(\mathfrak n^+=\bigoplus_{\alpha\in\Phi^+}\mathfrak g_\alpha\)、\(\mathfrak b=\mathfrak h\oplus\mathfrak n^+\)，给定 \(\lambda\in\mathfrak h^*\)。令 \(\mathbb C_\lambda\) 为一维 \(\mathfrak b\) 模，其中每个 \(h\in\mathfrak h\) 按 \(\lambda(h)\) 作用，\(\mathfrak n^+\) 作用为零。诱导模
	\[
		M(\lambda)=U(\mathfrak g)\otimes_{U(\mathfrak b)}\mathbb C_\lambda
	\]
	称为最高权 \(\lambda\) 的Verma 模，记 \(v_\lambda=1\otimes1\)。
\end{definition}

这个一维作用保持括号，因为 \([\mathfrak b,\mathfrak b]=\mathfrak n^+\) 被送到零。等价地，\(M(\lambda)\) 是 \(U(\mathfrak g)\) 模掉由全部 \(x\in\mathfrak n^+\) 及 \(h-\lambda(h)\)、\(h\in\mathfrak h\) 生成的左理想。不能把它说成双边理想商代数：这里构造的是一个左模。

\begin{proposition}\label{b5:prop:verma-pbw-universal}
	设 \(\mathfrak g\) 为有限维复半单 Lie 代数，\(\mathfrak h\) 为其 Cartan 子代数，\(\Phi^+\) 为相应根系的一组正根，\(\lambda\in\mathfrak h^*\)。令 \(\mathfrak n^\pm=\bigoplus_{\alpha\in\Phi^+}\mathfrak g_{\pm\alpha}\)，记相应 Verma 模为 \(M(\lambda)\)、其标准最高权生成元为 \(v_\lambda=1\otimes1\)。映射
	\[
		U(\mathfrak n^-)\longrightarrow M(\lambda),\qquad u\longmapsto uv_\lambda
	\]
	是向量空间同构。特别地 \(v_\lambda\ne0\)，\(M(\lambda)_\lambda=\mathbb Cv_\lambda\)，且其权空间有限维。 对任意 \(\mathfrak g\) 模 \(V\)，若 \(v\in V\) 满足 \(hv=\lambda(h)v\)、\(\mathfrak n^+v=0\)，则存在唯一模同态 \(M(\lambda)\to V\) 将 \(v_\lambda\) 送到 \(v\)。
\end{proposition}
\begin{proof}
	PBW 给出乘法同构
	\[
		U(\mathfrak n^-)\otimes_{\mathbb C}U(\mathfrak b) \longrightarrow U(\mathfrak g)
	\]
	作为向量空间以及右 \(U(\mathfrak b)\) 模的同构。再在右边张量 \(\mathbb C_\lambda\)，得到所需同构；常数单项式对应 \(v_\lambda\)，所以它非零。各 PBW 单项式的权按\cref{b5:prop:highest-weight-support}的计算给出，得到其余权空间结论。

	映射 \(u\otimes1\mapsto uv\) 良定，因为 \(v\) 满足 \(\mathbb C_\lambda\) 的全部 \(U(\mathfrak b)\) 关系；它显然保持左作用。反之，\(v_\lambda\) 生成 \(M(\lambda)\)，所以它的像决定整个同态。若 \(v\) 生成 \(V\)，这个同态满射。
\end{proof}

取 \(\mathfrak g=\mathfrak{sl}_3(\mathbb C)\)，以迹零对角矩阵为 Cartan 子代数，严格上三角矩阵给出正根向量。记 \(\varepsilon_j\) 为第 \(j\) 个对角坐标，\(\alpha_1=\varepsilon_1-\varepsilon_2\)、\(\alpha_2=\varepsilon_2-\varepsilon_3\)，并取
\[
	f_1=E_{21},\qquad f_2=E_{32},\qquad f_{12}=[f_2,f_1]=E_{31}.
\]
基本权 \(\omega_1=\varepsilon_1\) 在 \(h_1=E_{11}-E_{22}\)、\(h_2=E_{22}-E_{33}\) 上分别取值 \(1,0\)。按 \(f_1<f_{12}<f_2\) 排序，\(M(\omega_1)\) 的 PBW 基为 \(f_1^a f_{12}^b f_2^c v_{\omega_1}\)，其中 \(a,b,c\in\mathbb Z_{\ge0}\)，其权为
\[
	\omega_1-(a+b)\alpha_1-(b+c)\alpha_2.
\]
要得到权 \(\omega_1-\alpha_1-\alpha_2\)，必须有 \(a+b=b+c=1\)，只有 \((a,b,c)=(0,1,0),(1,0,1)\) 两种可能。因此
\[
	M(\omega_1)_{\omega_1-\alpha_1-\alpha_2} =\mathbb C f_{12}v_{\omega_1} \oplus\mathbb C f_1f_2v_{\omega_1}.
\]

\ProseParagraphBreak
在自然模 \(V=\mathbb C^3\) 的标准基 \(v_1,v_2,v_3\) 中，\(v_1\) 的权为 \(\omega_1\)，被全部正根向量零化，且 \(f_1v_1=v_2\)、\(f_2v_2=v_3\)，所以它生成 \(V\)。泛性质给出满射 \(M(\omega_1)\to V\)，将 \(v_{\omega_1}\) 送到 \(v_1\)。这个自然模简单：非零子模在对角矩阵作用下按三个互异的权分解，故含有某个 \(v_j\)；矩阵单位 \(E_{ij}\)、\(i\ne j\)，再由它生成全部标准基。目标权 \(\omega_1-\alpha_1-\alpha_2=\varepsilon_3\) 的权空间因而为 \(\mathbb Cv_3\)，而两个 PBW 方向的像为
\[
	f_{12}v_1=v_3,\qquad f_1f_2v_1=0,\qquad f_2f_1v_1=v_3.
\]
所以 Verma 模中该权的重数为 \(2\)，在这个简单商中却为 \(1\)：PBW 基在取商后可能消失或出现线性关系。并且 \(f_2v_1=0\) 只规定最高权生成元上的作用，\(f_2v_2=v_3\ne0\) 表明 \(f_2\) 在整个模上并不按零作用。

\ProseParagraphBreak
例如对 \(\mathfrak{sl}_2\)，写 \(\lambda(h)=z\in\mathbb C\)。则 \(M(z)\) 的基为 \(v,fv,f^2v,\ldots\)，作用为
\begin{equation}\label{b5:eq:sl2-verma-actions}
	h f^jv=(z-2j)f^jv,\qquad e f^jv=j(z-j+1)f^{j-1}v.
\end{equation}
每个权空间一维，但模本身无限维；“每个权空间有限维”与“只有有限多个权”是不同的条件。

\subsection{唯一简单商}
\label{b5:subsec:1803:02}

\begin{theorem}\label{b5:thm:verma-simple-quotient}
	设 \(\mathfrak g\) 为有限维复半单 Lie 代数，\(\mathfrak h\) 为其 Cartan 子代数，固定相应根系的一组正根，\(\lambda\in\mathfrak h^*\)。相对于这份选择的 Verma 模 \(M(\lambda)\) 有唯一极大真子模 \(N(\lambda)\)，因此有唯一简单商
	\[
		L(\lambda)=M(\lambda)/N(\lambda).
	\]
	每个非零最高权为 \(\lambda\) 的模都有一个满射到 \(L(\lambda)\)。所有简单最高权模由 \(\lambda\in\mathfrak h^*\) 唯一参数化。
\end{theorem}
\symidx{\(L(\lambda)\)：Verma 模的唯一简单商}
\begin{proof}
	由\cref{b5:prop:highest-weight-support}，任意子模都按权分解。一个真子模不可能含有 \(v_\lambda\)，所以它的 \(\lambda\) 权空间为零。令 \(N(\lambda)\) 为所有真子模之和；这里的子模之和按代数意义定义，其元素均为有限多个真子模中元素之和。每个真子模的 \(\lambda\) 权分量都为零，故 \(N(\lambda)_\lambda=0\)。因此 \(N(\lambda)\ne M(\lambda)\)，却包含每个真子模，它就是唯一极大真子模。商模非零且没有非零真子模，所以简单。

	若 \(V\ne0\) 为最高权 \(\lambda\) 的模，由泛性质有满射 \(M(\lambda)\to V\)。其核为真子模，故包含于 \(N(\lambda)\)，因此满射 \(M(\lambda)\to L(\lambda)\) 经 \(V\) 分解。若 \(V\) 简单，这个非零满射就是同构。不同最高权的简单模不能同构，因为一个最高权模的最高权生成元只有一个权，已在上一节证明。
\end{proof}

由此，上面算出的三维自然模同构于 \(L(\omega_1)\)。一般模的真子模之和却可能等于全模。例如取\cref{b5:prop:sl2-model}中的二维简单 \(\mathfrak{sl}_2(\mathbb C)\) 模 \(L_1\)，在 \(V=L_1\oplus L_1\) 中，
\[
	U_1=L_1\oplus0,\qquad U_2=0\oplus L_1
\]
都是非零真子模，却有 \(U_1+U_2=V\)。它们与 \(V\) 的二维权 \(1\) 空间相交于不同的直线。Verma 模的每个真子模则都缺失同一条最高权线 \(\mathbb Cv_\lambda\)；子模按权分解，故有限求和不能产生任何非零的 \(\lambda\) 权分量。这才使所有真子模之和仍然为真，而不只是因为模有一个循环生成元。

\ProseParagraphBreak
作为一个已经算过的例子，\cref{b5:thm:sl2-verma-submodules}给出
\[
	0\longrightarrow M(-m-2)\longrightarrow M(m)\longrightarrow L_m\longrightarrow0 \qquad(m\in\mathbb Z_{\ge0}),
\]
且此列不分裂。左端最高向量映到 \(f^{m+1}v\)，它是一个较低权的最高权向量；右端则是唯一有限维简单商。这个例子把“模中有另一个最高权向量”与“另有一个最高权生成元”的差别具体表现出来。

\subsection{有限维性的支配整性条件}
\label{b5:subsec:1803:03}

\begin{definition}\label{b5:def:dominant-integral-weight}
	设 \(\mathfrak g\) 为有限维复半单 Lie 代数，\(\mathfrak h\) 为其 Cartan 子代数，\(\alpha_1,\ldots,\alpha_\ell\) 为相应根系中选定的简单根，\(h_1,\ldots,h_\ell\in\mathfrak h\) 为对应余根。给定 \(\lambda\in\mathfrak h^*\)。若满足
	\[
		\lambda(h_i)\in\mathbb Z_{\ge0}\qquad(1\le i\le\ell),
	\]
	就称 \(\lambda\) 为支配整权。它们的集合记为 \(P^+\)。
\end{definition}

\begin{proposition}\label{b5:prop:finite-highest-dominant}
	设 \(\mathfrak g\) 为有限维复半单 Lie 代数，固定 Cartan 子代数 \(\mathfrak h\)、正根及标准简单根三元组 \(e_i,f_i,h_i\)（\(1\le i\le\ell\)）。令 \(P^+\) 为相对于此选择的支配整权集。若 \(V\ne0\) 为有限维复 \(\mathfrak g\) 最高权模，\(v\) 为最高权 \(\lambda\) 的生成元，则 \(\lambda\in P^+\)。令 \(m_i=\lambda(h_i)\)，则对每个 \(1\le i\le\ell\)，有
	\[
		f_i^{m_i+1}v=0.
	\]
\end{proposition}
\begin{proof}
	限制到第 \(i\) 个 \(\mathfrak{sl}_2\) 子代数。有限维完全可约性把 \(V\) 分成若干 \(L_m\)（\(m\in\mathbb Z_{\ge0}\)）的直和。\(v\) 被 \(e_i\) 零化且为 \(h_i\) 的特征向量，所以它在各直和项中的非零投影都是最高权为 \(\lambda(h_i)\) 的向量。由 \(\mathfrak{sl}_2\) 分类，这个数必须是非负整数 \(m_i\)，且这些投影都被 \(f_i^{m_i+1}\) 零化。把投影相加即得结论。
\end{proof}

到这里只证明了必要性。给定支配整权，是否存在有限维表示，不能从“这些关系看起来有限”直接推出：不同负根方向可以产生任意长的 PBW 单项式。下一节将先把简单根方向的幂关系加入 Verma 模，再证明这些关系控制了整个权集。
